% Guide des modules syllabus : mod_syllabus, mod_syllabusstaff, mod_syllabusstudent.
% Compilation : latexmk -pdf guide-syllabus.tex
% Captures : SYLLABUS=syllabus.json node captures.js site-x.json sy-... (voir captures.js).
\documentclass[11pt,a4paper]{article}

\newcommand{\titreguide}{Syllabus et évaluation des modules}
\newcommand{\versionguide}{1.1}
\newcommand{\dateguide}{novembre 2026}
\newcommand{\versionep}{1.1.0}
\newcommand{\versionsynthese}{1.0.0}

\usepackage{guide-commun}
\setcounter{tocdepth}{2}
\newlist{etapes}{enumerate}{1}
\setlist[etapes]{label=\textbf{\arabic*.}, leftmargin=*, itemsep=3pt}

% La page de garde commune cite mod_ep et mod_epsynthesis : celle-ci cite les trois modules.
\renewcommand{\pagedegarde}[2]{%
  \begin{titlepage}
    \centering
    \vspace*{3cm}
    {\Huge\bfseries \titreguide\par}
    \vspace{0.6cm}
    {\large #1\par}
    \vspace{1.5cm}
    \setlength{\fboxsep}{0pt}
    \fcolorbox{cadre}{white}{\includegraphics[width=0.85\linewidth]{captures/#2}}
    \vfill
    {\large Sébastien Lefebvre\par}
    \vspace{0.3cm}
    {\large Version \versionguide{} -- \dateguide\par}
    {\small Modules couverts : \texttt{mod\_syllabus}, \texttt{mod\_syllabusstaff},
     \texttt{mod\_syllabusstudent} \versionep\par}
  \end{titlepage}}

\begin{document}

\pagedegarde{La fiche de chaque module d'enseignement, sa relecture par la DEVE, et son
évaluation par les étudiants}{sy-etu-mes-modules}

\tableofcontents

\section*{Historique des versions}
\addcontentsline{toc}{section}{Historique des versions}
\begin{center}
\begin{tabularx}{\linewidth}{@{}l l X@{}}
\toprule
\textbf{Version} & \textbf{Date} & \textbf{Modifications} \\
\midrule
1.0 & novembre 2026 & Première version. Captures prises sur le banc : 12 modules, 24
enseignants, 200 étudiants, après rédaction, relecture et évaluation de tous les modules. \\
1.1 & novembre 2026 & Référentiel de compétences : modules et cours liés aux compétences, avec
un niveau ; matrice compétences $\times$ cours. \\
\bottomrule
\end{tabularx}
\end{center}

\newpage
% ===========================================================================
\part{Principes}\label{part:principes}
% ===========================================================================

\section{Trois activités}

\begin{description}
  \item[Syllabus du module] (\texttt{mod\_syllabus}) une activité dans le cours de chaque module
    d'enseignement. Elle porte la \textbf{fiche} du module (au modèle ECTS) et son
    \textbf{évaluation} par les étudiants du cours.
  \item[Suivi des modules] (\texttt{mod\_syllabusstaff}) une activité dans un cours d'équipe.
    Chaque enseignant y retrouve les modules dont il est responsable, dans tous les cours. La
    DEVE y voit tous les modules, relit les fiches, fixe les évaluations, compare les résultats
    et tient le \textbf{référentiel commun} (types de module, questions communes, réglages).
  \item[Mes modules] (\texttt{mod\_syllabusstudent}) une activité dans le cours de la
    promotion. Chaque étudiant y voit les modules des cours où il est inscrit, leur fiche, et
    les évaluations qu'il a à faire.
\end{description}

Les deux activités de suivi se limitent, au choix, à une catégorie de cours (sous-catégories
comprises) ou couvrent tout le site.

\section{Compétences}\label{sec:competences}

Le \textbf{référentiel de compétences} est celui de Moodle (\chemin{Administration du site >
Compétences > Référentiels de compétences}) : un ou plusieurs référentiels, chacun en arbre
(domaines, puis compétences), créés à la main ou importés d'un fichier CSV. Les modules et les
cours s'y rattachent :
\begin{itemize}
  \item un \textbf{module} vise des compétences, avec un \textbf{niveau} : \statut{Introduite},
    \statut{Approfondie} ou \statut{Maîtrisée} ;
  \item un \textbf{cours} peut viser des compétences sans passer par un module (page
    \og Compétences\fg{} du cours, dans Moodle).
\end{itemize}
Les liens sont ceux de Moodle (compétence liée au cours, puis à l'activité) : ils servent aussi
aux plans de formation et aux preuves de Moodle. Le niveau est propre au module.

\section{Qui fait quoi}

\begin{center}
\begin{tabularx}{\linewidth}{@{}l X@{}}
\toprule
\textbf{Rôle} & \textbf{Ce qu'il fait} \\
\midrule
Responsable du module & Enseignant désigné : rédige la fiche, pose ses questions
d'évaluation, fixe la fenêtre d'évaluation, voit les résultats (après la fermeture, sauf
réglage contraire). \\
DEVE & Tout : désigne les responsables, relit et publie les fiches, fixe les fenêtres en lot,
définit les types de module et les questions communes, voit tous les résultats. \\
Étudiant & Lit les fiches, évalue chaque module pendant sa fenêtre, une fois, anonymement. \\
\bottomrule
\end{tabularx}
\end{center}

\section{La fiche et ses versions}

Une fiche passe par des \textbf{versions} :
\begin{center}
\begin{tabularx}{\linewidth}{@{}l X@{}}
\toprule
\textbf{État} & \textbf{Signification} \\
\midrule
\statut{Brouillon} & En cours de rédaction ; les étudiants voient la version publiée. \\
\statut{Soumise à la DEVE} & Attend la relecture ; ne se modifie plus. \\
\statut{Renvoyée par la DEVE} & À reprendre, avec le commentaire de la DEVE. \\
\statut{Publiée} & Celle que lisent les étudiants. \\
\statut{Refusée} & Écartée par la DEVE. \\
\statut{Remplacée} & Ancienne version publiée, gardée dans l'historique. \\
\bottomrule
\end{tabularx}
\end{center}

Si la DEVE a désactivé la relecture (référentiel), la fiche est publiée dès sa soumission.

\section{L'évaluation}\label{sec:principe-evaluation}

\begin{itemize}
  \item Elle s'ouvre à une \textbf{date} (et se ferme à une autre, facultative). À l'ouverture,
    les étudiants du cours sont prévenus par un message.
  \item Les étudiants répondent d'abord aux \textbf{questions de la DEVE} : celles posées à
    tous les modules, puis celles posées aux modules de leur type ; ensuite aux
    \textbf{questions de l'enseignant}.
  \item Formes de réponse : échelle de 1 (pas du tout d'accord) à 5 (tout à fait d'accord),
    oui ou non, choix dans une liste, réponse libre ; chaque question peut être obligatoire.
  \item À l'ouverture, les questions sont \textbf{figées} : celles de la DEVE sont copiées dans
    le module. Les modifier ensuite ne change pas une évaluation déjà ouverte, et les résultats
    portent toujours sur les questions effectivement posées.
  \item Les réponses sont \textbf{anonymes} : la réponse ne porte ni le nom de l'étudiant ni
    l'heure (seulement le jour) ; seul le fait qu'il a répondu est noté, à part, pour qu'il ne
    réponde qu'une fois.
  \item Les résultats ne s'affichent qu'à partir d'un \textbf{nombre minimal de réponses}
    (réglage de la DEVE, 3 par défaut), pour qu'on ne puisse pas deviner qui a répondu quoi.
\end{itemize}

\begin{attention}
L'anonymat vaut pour toutes les pages et tous les exports. Un administrateur qui lirait
directement la base de données pourrait, en rapprochant l'ordre des enregistrements, deviner
approximativement l'auteur d'une réponse : l'accès à la base doit rester restreint, comme pour
toute donnée sensible.
\end{attention}

\newpage
% ===========================================================================
\part{Étudiant}\label{part:etudiant}
% ===========================================================================

\section{Mes modules}

Dans le cours de votre promotion, ouvrez l'activité \og Mes modules\fg{}. En tête, les
\textbf{évaluations à faire}, la plus urgente d'abord ; en dessous, tous vos modules : type,
ECTS, période, responsables, et où vous en êtes de l'évaluation (\statut{Évalué}, ouverture à
venir, close).

\capture{sy-etu-mes-modules}{Mes modules : évaluations à faire, puis tous les modules.}

Le nom d'un module ouvre sa \textbf{fiche} : caractéristiques, objectifs, acquis
d'apprentissage, contenu, méthodes, modalités d'évaluation, programme des séances et volume
horaire, et les \textbf{compétences visées} avec leur niveau. Elle se télécharge en PDF.

\capture{sy-etu-fiche}{Fiche d'un module, avec l'invitation à l'évaluer.}

\section{Évaluer un module}

\begin{etapes}
  \item Cliquez sur \bouton{Évaluer} (dans \og Mes modules\fg{}) ou \bouton{Évaluer ce module}
    (sur la fiche).
  \item Répondez aux questions de la DEVE, puis à celles de l'enseignant. Les questions
    marquées d'un astérisque sont obligatoires ; aucune réponse n'est cochée d'avance.
  \item Cliquez sur \bouton{Envoyer mon évaluation}.
\end{etapes}

\capture{sy-etu-evaluer}{Formulaire d'évaluation.}

Vous ne pouvez répondre qu'une fois, pendant la fenêtre d'évaluation. Vos réponses sont
anonymes (section~\ref{sec:principe-evaluation}).

\newpage
% ===========================================================================
\part{Enseignant responsable}\label{part:enseignant}
% ===========================================================================

\section{Mes modules}

L'activité \og Suivi des modules\fg{} du cours d'équipe liste les modules dont vous êtes
responsable : état de la fiche, de l'évaluation, nombre de réponses, et les liens vers la fiche,
l'évaluation et les résultats. Un bandeau signale les modules sans fiche publiée ou dont la
fiche a été renvoyée.

\capture{sy-ens-mes-modules}{Suivi des modules : mes modules.}

Sur la page du module, deux encadrés résument l'évaluation (fenêtre, réponses) et la rédaction
(version publiée, version en cours), avec les boutons utiles.

\capture{sy-ens-fiche}{Page du module vue par son responsable.}

\section{Rédiger la fiche}

\begin{etapes}
  \item \bouton{Rédiger la fiche} (ou \bouton{Mettre à jour la fiche}) : la version en cours
    reprend la version publiée.
  \item Renseignez les rubriques : objectifs, acquis d'apprentissage visés, prérequis, contenu,
    méthodes, modalités d'évaluation, seconde session, bibliographie ; la langue
    d'enseignement et les heures de travail personnel.
  \item Décrivez le \textbf{programme des séances} : intitulé, nature (cours magistral,
    travaux dirigés, travaux pratiques, clinique, terrain, projet, autre), durée et intervenant
    (un enseignant du cours, ou le nom d'un intervenant extérieur). Les séances ne sont pas
    datées : la fiche décrit le module, pas le calendrier de l'année.
  \item \bouton{Enregistrer le brouillon} autant que nécessaire, puis \bouton{Enregistrer et
    soumettre à la DEVE} (ou \bouton{Enregistrer et publier} si la relecture est désactivée).
\end{etapes}

\capture{sy-ens-redaction}{Rédaction de la fiche.}

\capture{sy-ens-seances}{Programme des séances.}

Les caractéristiques du module (code, ECTS, période, année d'étude, type de module) se règlent
dans les \chemin{Paramètres} de l'activité. Les \chemin{Responsables} se choisissent parmi les
enseignants du cours ; celui qui ajoute l'activité en devient responsable.

\capture[0.85\linewidth]{sy-ens-responsables}{Responsables du module.}

\section{Lier le module aux compétences}

\begin{etapes}
  \item \chemin{Compétences} (onglets du module, ou menu \og Plus\fg{}) : la page rappelle les
    compétences déjà visées.
  \item Choisissez le \textbf{référentiel} : ses compétences s'affichent en arbre.
  \item Pour chaque compétence que le module vise, choisissez le niveau ; laissez
    \og Non visée\fg{} pour les autres.
  \item \bouton{Enregistrer les modifications}.
\end{etapes}

\capture{sy-ens-competences}{Compétences visées par le module et liaison au référentiel.}

Une compétence retirée du module reste liée au cours : celui-ci peut la viser par ailleurs.
Lier des compétences demande, en plus d'être responsable, le droit de gérer les compétences du
cours (enseignant éditeur).

\section{Préparer l'évaluation}

\chemin{Évaluation du module} réunit :
\begin{itemize}
  \item la \textbf{fenêtre} : ouverture et, au besoin, fermeture ;
  \item les \textbf{questions de la DEVE} qui s'appliqueront à votre module (pour information) ;
  \item \textbf{vos questions} : ajoutez-les, modifiez-les, réordonnez-les jusqu'à l'ouverture.
\end{itemize}

\capture{sy-ens-evaluation}{Évaluation du module ouverte : fenêtre et questions figées.}

\section{Lire les résultats}

\chemin{Résultats de l'évaluation} donne le nombre de réponses sur le nombre d'étudiants, et
pour chaque question la moyenne (échelle), la répartition des réponses, ou la liste des
réponses libres (dans l'ordre alphabétique, rien ne rappelant l'ordre des réponses). Les
résultats s'exportent en CSV. Ils sont visibles après la fermeture de l'évaluation, sauf si la
DEVE les a ouverts dès l'ouverture, et seulement au-delà du nombre minimal de réponses.

\capture{sy-ens-resultats}{Résultats de l'évaluation d'un module.}

\newpage
% ===========================================================================
\part{DEVE}\label{part:deve}
% ===========================================================================

\section{Mise en place}

\begin{etapes}
  \item Dans un cours d'équipe, ajoutez une activité \og Suivi des modules\fg{} (périmètre :
    la catégorie des cours de modules, ou tout le site). Donnez-vous le rôle
    \og gestionnaire\fg{} dans ce cours.
  \item Dans \chemin{Référentiel commun} : réglages, types de module, questions communes.
  \item Dans chaque cours de module : une activité \og Syllabus du module\fg{} (ajoutée par
    l'enseignant, ou par vous) ; vous devez y avoir le rôle \og gestionnaire\fg{} (par exemple
    au niveau de la catégorie) pour relire et fixer les fenêtres.
  \item Dans le cours de chaque promotion : une activité \og Mes modules\fg{}.
\end{etapes}

\section{Référentiel commun}

Il vaut pour tous les modules du site.
\begin{description}
  \item[Réglages] relecture des fiches par la DEVE avant publication ; nombre minimal de
    réponses pour afficher des résultats ; résultats visibles des responsables après la
    fermeture ou dès l'ouverture.
  \item[Types de module] (fondamental, clinique, optionnel\dots) : chaque module en porte un
    (dans ses paramètres) ; un type porté par des modules ne se supprime pas.
  \item[Questions communes] posées à tous les modules, ou aux modules d'un type ; forme de
    réponse, obligation, activation, ordre.
\end{description}

\capture{sy-deve-referentiel}{Référentiel commun : réglages, types, questions.}

\section{Tous les modules}

\chemin{Tous les modules} liste les modules du périmètre, filtrables par nom ou code, type,
état de la fiche et état de l'évaluation. Cochez des modules, saisissez une ouverture (et une
fermeture), puis \bouton{Appliquer} : la même fenêtre est fixée pour tous, là où vous avez le
rôle de la DEVE (les autres sont signalés). Laisser l'ouverture vide retire l'évaluation.

\capture{sy-deve-tous}{Tous les modules, avec la fenêtre d'évaluation en lot.}

\section{Relire les fiches}

\chemin{Fiches à relire} liste les fiches soumises. \bouton{Relire} ouvre la version proposée,
avec le formulaire de décision : \bouton{Publier}, \bouton{Renvoyer aux responsables} ou
\bouton{Refuser} (un commentaire est alors obligatoire). Les responsables sont prévenus.

\capture{sy-deve-a-relire}{Fiches à relire.}

\capture{sy-deve-relecture}{Relecture d'une fiche soumise.}

\section{Matrice des compétences}

\chemin{Compétences} (suivi des modules) croise un référentiel avec les cours du périmètre :
une ligne par compétence, une colonne par cours. Une case porte le niveau donné par le module
(I, A, M), \og $\checkmark$\fg{} pour un lien au module sans niveau, \og \textbullet\fg{} pour un lien au cours
seulement. La dernière colonne compte les cours qui visent la compétence ; une compétence
qu'aucun cours ne vise est surlignée. Filtre par année d'étude, export CSV. Les enseignants y
ont accès aussi.

\capture{sy-deve-competences}{Matrice compétences $\times$ cours, avec la couverture du référentiel.}

L'en-tête d'une colonne ouvre la page \og Compétences\fg{} du module (ou, pour un cours sans
module, la page des compétences du cours de Moodle).

\section{Résultats comparés}

\chemin{Résultats des évaluations} donne, module par module, le nombre de réponses et la
moyenne de chaque question commune à échelle : les mêmes questions, posées partout, se
comparent d'un module à l'autre. Les modules sous le seuil d'anonymat n'ont pas de moyenne.
Export CSV.

\capture{sy-deve-resultats}{Résultats comparés des évaluations closes.}

\newpage
% ===========================================================================
\part{Administrateur du site}\label{part:admin}
% ===========================================================================

\section{Installation}

Copier \texttt{mod/syllabus}, \texttt{mod/syllabusstaff} et \texttt{mod/syllabusstudent} dans
\texttt{<moodle>/mod/}, puis \texttt{php admin/cli/upgrade.php}. Moodle 4.3 ou plus récent.
Les deux activités de suivi dépendent de \texttt{mod\_syllabus}. Une tâche planifiée,
\texttt{mod\_syllabus\textbackslash task\textbackslash open\_evaluations} (tous les quarts
d'heure), fige les questions des évaluations qui viennent d'ouvrir et prévient les étudiants :
le \emph{cron} doit tourner.

\begin{center}
\small
\begin{tabularx}{\linewidth}{@{}l X@{}}
\toprule
\textbf{Capacité} & \textbf{Pour qui (par défaut)} \\
\midrule
\texttt{mod/syllabus:view} & Tous : lire la fiche publiée \\
\texttt{mod/syllabus:edit} & Enseignants : pouvoir être responsable \\
\texttt{mod/syllabus:manage} & Gestionnaire (DEVE) : responsables, relecture, fenêtre,
résultats \\
\texttt{mod/syllabus:evaluate} & Étudiants : évaluer \\
\texttt{mod/syllabusstaff:view} & Enseignants : leurs modules \\
\texttt{mod/syllabusstaff:manageall} & Gestionnaire (DEVE) : tous les modules, référentiel \\
\texttt{mod/syllabusstudent:view} & Étudiants : leurs modules \\
\texttt{moodle/competency:coursecompetencymanage} & Enseignants éditeurs (Moodle) : lier des
compétences au cours et au module \\
\bottomrule
\end{tabularx}
\end{center}

Les compétences doivent être activées sur le site (\chemin{Administration du site >
Fonctions avancées}, activé par défaut) ; les référentiels se gèrent dans \chemin{Administration
du site > Compétences}, avec \texttt{moodle/competency:competencymanage}.

La sauvegarde Moodle du cours comprend la fiche, ses versions et les questions du module ;
avec les données des utilisateurs, les responsables et l'évaluation (réponses anonymes et
participations), et les liens aux compétences avec leur niveau (retrouvés par les identifiants
du référentiel et de la compétence sur un autre site). La réinitialisation du cours peut effacer les réponses : l'évaluation est
alors de nouveau modifiable. Le référentiel commun n'appartient à aucun cours : il n'est pas
sauvegardé avec eux.

\newpage
% ===========================================================================
\part{Questions fréquentes}\label{part:faq}
% ===========================================================================

\begin{description}
  \item[Un enseignant du cours ne peut pas rédiger la fiche.] Il n'est pas désigné responsable
    (\chemin{Responsables}). Tant qu'un module n'a aucun responsable, tout enseignant du cours
    peut le rédiger.
  \item[Je ne peux plus modifier mes questions.] L'évaluation est ouverte : les questions sont
    figées pour que toutes les réponses portent sur les mêmes questions.
  \item[Je ne vois pas les résultats.] Ils s'affichent après la fermeture (sauf réglage de la
    DEVE) et à partir du nombre minimal de réponses.
  \item[Un étudiant ne voit pas un module dans « Mes modules ».] Il n'est pas inscrit au cours
    du module, l'activité est cachée, ou le cours est hors du périmètre de l'activité.
  \item[La DEVE a modifié une question commune, mais l'ancienne est posée.] L'évaluation de ce
    module était déjà ouverte : elle garde ses questions. La nouvelle vaudra pour les
    évaluations ouvertes ensuite.
  \item[La page « Compétences » du module ne propose pas de lier.] Il faut le droit de gérer les
    compétences du cours (enseignant éditeur) ; un enseignant non éditeur, même responsable,
    voit les compétences sans pouvoir les modifier.
  \item[Aucun référentiel n'est proposé.] Aucun référentiel visible n'existe : l'administrateur
    ou la DEVE en crée ou en importe un dans \chemin{Administration du site > Compétences}.
  \item[Peut-on savoir qui a répondu quoi ?] Non, par aucune page ni aucun export ; seul le fait
    d'avoir répondu est noté (voir l'encadré de la section~\ref{sec:principe-evaluation}).
\end{description}

\end{document}
